% Источники: 07_Программные_требования/DOC-06.3_Ранцевые_криптосистемы.md
% Генератор: docs/build_pdf.py
% Шаблон печатного представления. Содержимое подставляет docs/build_pdf.py.
\documentclass[12pt,a4paper]{article}
\usepackage{iftex}
\ifPDFTeX
  \usepackage[utf8]{inputenc}
  \usepackage[T2A]{fontenc}
  \usepackage{newtxtext,newtxmath}
\else
  \usepackage{fontspec}
  \IfFontExistsTF{Times New Roman}
    {\setmainfont{Times New Roman}}
    {\setmainfont{Libertinus Serif}}
  \IfFontExistsTF{Arial}{\setsansfont{Arial}}{\setsansfont{Latin Modern Sans}}
  \IfFontExistsTF{Menlo}{\setmonofont{Menlo}[Scale=MatchLowercase]}
    {\setmonofont{Latin Modern Mono}[Scale=MatchLowercase]}
\fi
\usepackage[russian]{babel}
\usepackage{geometry}
\geometry{left=25mm,right=15mm,top=20mm,bottom=20mm}
\usepackage{setspace}
\onehalfspacing
\usepackage{indentfirst}
\setlength{\parindent}{1.25cm}
\setlength{\parskip}{2pt}
\setlength{\emergencystretch}{3em}
\usepackage{graphicx,booktabs,longtable,array,tabularx,calc,etoolbox}
\usepackage{enumitem}
\setlist[itemize]{leftmargin=1.25cm,itemsep=2pt,topsep=4pt}
\setlist[enumerate]{leftmargin=1.25cm,itemsep=2pt,topsep=4pt}
\newcommand{\tightlist}{\setlength{\itemsep}{0pt}\setlength{\parskip}{0pt}}
\AtBeginEnvironment{longtable}{\small\singlespacing}
\setlength{\LTpre}{8pt}
\setlength{\LTpost}{8pt}
\setlength{\tabcolsep}{5pt}
\setlength{\LTleft}{0pt}
\setlength{\LTright}{0pt}
\usepackage{xurl}
\usepackage[unicode,colorlinks=true,linkcolor=black,urlcolor=blue]{hyperref}
\usepackage{bookmark}
\hypersetup{pdftitle={Ранцевые криптосистемы},pdfauthor={Георгиев Д.Д.},pdfsubject={Криптотренажёр mephi-12}}
\urlstyle{same}
\pagestyle{plain}
\setcounter{secnumdepth}{3}
\setcounter{tocdepth}{2}

\begin{document}
\hypersetup{pageanchor=false}
\begin{titlepage}
  \thispagestyle{empty}
  \centering
  {\large НИЯУ МИФИ\par}
  \vspace*{3.5cm}
  {\LARGE\bfseries Ранцевые криптосистемы\par}
  \vspace{0.5cm}
  {\large DOC-06.3\par}
  \vspace{0.6cm}
  {\large Платформа для изучения и практики курса по криптографии\par}
  {\large Криптотренажёр mephi-12\par}
  \vspace{0.7cm}
  {\large Версия 0.9\par}
  {\large Редакция от 20.11.2026\par}
  \vspace{1cm}
  \begin{tabularx}{\textwidth}{@{}p{3.4cm}Xp{4cm}@{}}
    \toprule
    Роль & ФИО & Подпись / Дата \\
    \midrule
    Разработал & Георгиев Д.Д. & \\
    \addlinespace[6pt]
    Согласовано & & \\
    \addlinespace[6pt]
    Утверждено & & \\
    \bottomrule
  \end{tabularx}
  \vfill
  {\large Москва, 2026\par}
\end{titlepage}
\hypersetup{pageanchor=true}

\section*{Лист регистрации изменений}
\addcontentsline{toc}{section}{Лист регистрации изменений}
\begingroup\small\singlespacing
\noindent
\begin{tabularx}{\textwidth}{@{}p{1.2cm}p{2cm}p{2.5cm}Xp{1.8cm}@{}}
\toprule
Версия & Дата & Автор & Содержание изменения & Утвердил \\
\midrule
0.9 & 20.11.2026 & Анохин Г.А. & Уточнены критерии выбора варианта, проверки решения и представления результата. & \\
\bottomrule
\end{tabularx}
\endgroup
\clearpage
\begingroup\small\singlespacing
\tableofcontents
\endgroup
\clearpage

\hypertarget{doc-063-ранцевые-криптосистемы}{}

{\def\LTcaptype{none} % do not increment counter
\begin{longtable}[]{@{}
  >{\raggedright\arraybackslash}p{(\linewidth - 2\tabcolsep) * \real{0.3200}}
  >{\raggedright\arraybackslash}p{(\linewidth - 2\tabcolsep) * \real{0.6800}}@{}}
\toprule\noalign{}
\begin{minipage}[b]{\linewidth}\raggedright
Поле
\end{minipage} & \begin{minipage}[b]{\linewidth}\raggedright
Значение
\end{minipage} \\
\midrule\noalign{}
\endhead
\bottomrule\noalign{}
\endlastfoot
Версия & 0.9 \\
Дата редакции & 20.11.2026 \\
Автор предметных требований & Георгиев Д.Д. \\
Редактор версии 0.9 & Анохин Г.А. \\
Основания & Раздел {\href{../04\_Системные\_требования/DOC-05\_Системные\_требования.pdf\#doc-05-требования-к-ранцевым-задачам}{ранцевых задач}} DOC-05; {\href{DOC-06\_Программные\_требования.pdf\#doc-06-сводные-программные-требования}{DOC-06}} \\
Статус & Редакторская версия для согласования с владельцем подсистемы \\
\end{longtable}
}

\section{Назначение и состав}\label{1-ux43dux430ux437ux43dux430ux447ux435ux43dux438ux435-ux438-ux441ux43eux441ux442ux430ux432}

\hypertarget{1-назначение-и-состав}{}

Студент получает исходные данные учебного шифра, вычисляет недостающие величины и показывает, как из шифротекста получается исходное сообщение. В личной сессии платформа сохраняет выданное условие и проверяет решение студента. Для самоподготовки можно запросить отдельный пример с полным решением. В документе такая выдача называется редакционной генерацией. Её результат не входит в оценку сессии.

\subsection{Как устроена ранцевая задача}\label{ux43aux430ux43a-ux443ux441ux442ux440ux43eux435ux43dux430-ux440ux430ux43dux446ux435ux432ux430ux44f-ux437ux430ux434ux430ux447ux430}

\hypertarget{как-устроена-ранцевая-задача}{}

Ранец в этой задаче --- упорядоченный список положительных целых чисел, или весов. Позиции списка соответствуют битам сообщения. Бит равен нулю или единице. При шифровании складывают веса на тех позициях, где в сообщении стоит единица. Полученная сумма называется шифротекстом. Например, для весов 10, 15 и 1 сообщение из битов 1, 0, 1 даёт сумму 11.

\textbf{Лёгкий ранец} задаёт последовательность, в которой каждый элемент больше суммы предшествующих. Это свойство называют супервозрастанием. Когда известна сумма выбранных весов такого ранца, можно определить биты, рассматривая веса от большего к меньшему. Вариант «на степенях» строит последовательность из степеней одного основания. Вариант «на супервозрастающей последовательности» требует самого свойства без привязки к степеням.

\textbf{Тяжёлый ранец} получают из лёгкого. Каждый вес умножают на одно и то же число и берут остаток от деления на модуль. Модуль больше суммы весов лёгкого ранца. Множитель и модуль взаимно просты, то есть их единственный общий положительный делитель равен единице. Поэтому у множителя есть обратный по этому модулю --- число, произведение которого на множитель при делении на модуль даёт остаток один. При расшифровании обратный множитель позволяет перейти от шифротекста к сумме весов лёгкого ранца. Открытой частью ключа для шифрования служит тяжёлый ранец. Лёгкий ранец и параметры преобразования образуют закрытые данные для расшифрования. В учебном условии часть величин специально раскрывают студенту.

Далее описаны программные обязанности по {\href{../04\_Системные\_требования/DOC-05\_Системные\_требования.pdf\#doc-05-требования-к-ранцевым-задачам}{системным требованиям к ранцевым задачам}}. \textbf{Предметный обработчик} --- модуль, который выдаёт условие определённого типа и проверяет ответ по данным именно этой задачи. \textbf{Решатель} строит полное эталонное решение для самоподготовки.

{\def\LTcaptype{none} % do not increment counter
\begin{longtable}[]{@{}
  >{\raggedright\arraybackslash}p{(\linewidth - 2\tabcolsep) * \real{0.3200}}
  >{\raggedright\arraybackslash}p{(\linewidth - 2\tabcolsep) * \real{0.6800}}@{}}
\toprule\noalign{}
\begin{minipage}[b]{\linewidth}\raggedright
Логическая часть
\end{minipage} & \begin{minipage}[b]{\linewidth}\raggedright
Ответственность
\end{minipage} \\
\midrule\noalign{}
\endhead
\bottomrule\noalign{}
\endlastfoot
Участие в сессиях & По имени, под которым ранцевый модуль зарегистрирован в программе, направить ему запрос на выдачу или проверку задачи \\
Генерация условия & Построить условие выбранного варианта с частично скрытыми данными и сохранить его в составе сессии \\
Проверка решения & Проверить ответ по величинам криптосистемы и вернуть вердикт с диагностикой \\
Редакционные решатели & Построить эталонную задачу по варианту или с параметрами пользователя \\
Диспетчеризация & Выбрать решатель по варианту редакционного запроса \\
Хранение & Сохранять условия сгенерированных задач между запусками \\
\end{longtable}
}

В {\href{DOC-06.2\_Сессии\_задач\_и\_планировщик.pdf\#doc-062-контракты-сессий-задач}{контракте сессий}} заданы поля, типы и коды ошибок, в {\href{../09\_Архитектура\_и\_API/DOC-08\_Справочник\_API.pdf\#doc-08-справочник-api-сессии-задач}{разделе DOC-08}} описаны действия клиента. Для задач в сессии внешний тип --- \texttt{BACKPACK}. Поле \texttt{type} внутри решения указывает математический вариант выданного ранцевого задания.

\subsection{Условие и его состав}\label{ux443ux441ux43bux43eux432ux438ux435-ux438-ux435ux433ux43e-ux441ux43eux441ux442ux430ux432}

\hypertarget{условие-и-его-состав}{}

Условие ранцевой задачи хранится вместе с конкретной задачей сессии как объект JSON --- текстовое представление именованных полей со значениями. Оно содержит вариант, выданные фрагменты лёгкого ранца, тяжёлый ранец и данные, достаточные для восстановления остальных величин. После выдачи условие не формируется заново по шаблону: проверка использует его сохранённую копию, что соответствует {\href{DOC-06\_Программные\_требования.pdf\#swr-sys-003-состав-созданной-сессии}{составу созданной сессии}}.

Приоритет всех программных требований документа --- обязательный. Источники SR указывают на определения в DOC-05; сценарии приведены в {\hyperlink{7-трассировка-и-описательные-проверки}{таблице проверок}}.

\section{Участие в сессиях}\label{2-ux443ux447ux430ux441ux442ux438ux435-ux432-ux441ux435ux441ux441ux438ux44fux445}

\hypertarget{2-участие-в-сессиях}{}

\subsection{SWR-KNAP-001. Регистрация предметного обработчика}\label{swr-knap-001-ux440ux435ux433ux438ux441ux442ux440ux430ux446ux438ux44f-ux43fux440ux435ux434ux43cux435ux442ux43dux43eux433ux43e-ux43eux431ux440ux430ux431ux43eux442ux447ux438ux43aux430}

\hypertarget{swr-knap-001-регистрация-предметного-обработчика}{}

\hypertarget{swr-knap-001}{}

Управление сессиями получает задачи, заданные шаблоном, через предметных обработчиков. Квалификатор «Backpack» --- имя, под которым ранцевый модуль зарегистрирован для вызова из подсистемы сессий. По этому имени управление сессиями передаёт модулю формирование задания и проверку решения. Так состав набора задаётся без знания внутренних алгоритмов домена. Это применение {\href{DOC-06\_Программные\_требования.pdf\#swr-sys-004-контракт-предметного-обработчика}{общего контракта предметного обработчика}} для ранцевого типа.

\textbf{Требование.} Обработчик ранцевых задач должен быть доступен управлению сессиями под квалификатором «Backpack» и вызываться для задач, тип которых задан этим квалификатором в шаблоне.

Вызов генерации происходит при создании набора, вызов проверки --- после допуска ответа. Обработчик не получает и не изменяет правила срока, попытки и принадлежности: они остаются у управления сессиями. Порядок вызова определён в {\href{DOC-06.2\_Сессии\_задач\_и\_планировщик.pdf\#swr-ses-004-формирование-задач}{формировании задач}} и {\href{DOC-06.2\_Сессии\_задач\_и\_планировщик.pdf\#swr-ses-010-передача-профильной-проверке}{передаче профильной проверке}}.

Источники: {\href{../04\_Системные\_требования/DOC-05\_Системные\_требования.pdf\#sr-knap-001}{SR-KNAP-001}}; {\href{DOC-06\_Программные\_требования.pdf\#swr-sys-004}{SWR-SYS-004}}.

\subsection{SWR-KNAP-002. Случайный выбор варианта}\label{swr-knap-002-ux441ux43bux443ux447ux430ux439ux43dux44bux439-ux432ux44bux431ux43eux440-ux432ux430ux440ux438ux430ux43dux442ux430}

\hypertarget{swr-knap-002-случайный-выбор-варианта}{}

\hypertarget{swr-knap-002}{}

Ранцевая часть шаблона не различает учебные варианты: оба вида заданий обслуживаются одним квалификатором. Конкретный вариант каждого задания выбирается при генерации так, чтобы состав набора не предопределял тип всех ранцевых задач. Выбор не управляется параметрами пользователя и не зависит от предыдущих задач сессии.

\textbf{Требование.} При генерации задачи для сессии ПО должно выбирать вариант случайно из двух поддерживаемых.

Выбранный вариант фиксируется в сохранённом условии и далее используется проверкой без повторного выбора. Системное основание --- {\href{../04\_Системные\_требования/DOC-05\_Системные\_требования.pdf\#sr-knap-001-выбор-варианта-ранцевой-задачи}{выбор варианта ранцевой задачи}}.

Источник: {\href{../04\_Системные\_требования/DOC-05\_Системные\_требования.pdf\#sr-knap-001}{SR-KNAP-001}}.

\section{Генерация условия}\label{3-ux433ux435ux43dux435ux440ux430ux446ux438ux44f-ux443ux441ux43bux43eux432ux438ux44f}

\hypertarget{3-генерация-условия}{}

Студент получает условие с частью {\hyperlink{как-устроена-ранцевая-задача}{величин ранцевого шифра}}. Для последующей проверки у платформы остаются данные того же варианта. Правила ниже определяют размер сообщения, форму весов и состав выдачи. Размерность --- число весов в каждом ранце, равное числу битов сообщения.

\subsection{SWR-KNAP-003. Генерация сообщения}\label{swr-knap-003-ux433ux435ux43dux435ux440ux430ux446ux438ux44f-ux441ux43eux43eux431ux449ux435ux43dux438ux44f}

\hypertarget{swr-knap-003-генерация-сообщения}{}

\hypertarget{swr-knap-003}{}

Сообщение определяет, какие элементы тяжёлого ранца суммируются при шифровании, и входит как в условие, так и в проверяемый ответ. Учебное сообщение должно быть коротким для ручного решения и содержать достаточно нулей, чтобы расшифровка жадным алгоритмом была наглядной.

\textbf{Требование.} Сгенерированное сообщение должно быть случайной последовательностью логических значений длиной от 3 до 5 с долей ложных значений не менее трети.

Сгенерированное сообщение сохраняется в условии и сравнивается с переданным в ответе. Параметры длины относятся к задачам в сессиях; редакционные генерации принимают сообщение пользователя по {\hyperlink{swr-knap-013}{SWR-KNAP-013}}.

Источники: {\href{../04\_Системные\_требования/DOC-05\_Системные\_требования.pdf\#sr-knap-001}{SR-KNAP-001}}, {\href{../04\_Системные\_требования/DOC-05\_Системные\_требования.pdf\#sr-knap-002}{SR-KNAP-002}}.

\subsection{SWR-KNAP-004. Сокрытие параметров условия}\label{swr-knap-004-ux441ux43eux43aux440ux44bux442ux438ux435-ux43fux430ux440ux430ux43cux435ux442ux440ux43eux432-ux443ux441ux43bux43eux432ux438ux44f}

\hypertarget{swr-knap-004-сокрытие-параметров-условия}{}

\hypertarget{swr-knap-004}{}

Выдача полного набора данных превратила бы задание в чтение ответа. Поэтому условие скрывает большую часть лёгкого ранца, оставляя открытыми отдельные элементы, и может не содержать множитель. Степень сокрытия ограничена требованием восстанавливаемости: по условию должны однозначно достраиваться все величины криптосистемы.

\textbf{Требование.} В условии открытыми должны оставаться не более двух элементов лёгкого ранца; множитель преобразования может не включаться в условие.

Сокрытие выполняется при генерации и сохраняется вместе с условием: повторное чтение сессии выдаёт те же данные. Восстановление скрытого опирается на математические свойства последовательности варианта; полная проверка достроенных величин выполняется по {\hyperlink{swr-knap-009}{SWR-KNAP-009}}. Системное основание --- {\href{../04\_Системные\_требования/DOC-05\_Системные\_требования.pdf\#sr-knap-002-частично-скрытое-условие}{частично скрытое условие}}.

Источник: {\href{../04\_Системные\_требования/DOC-05\_Системные\_требования.pdf\#sr-knap-002}{SR-KNAP-002}}.

\subsection{SWR-KNAP-005. Супервозрастающий лёгкий ранец}\label{swr-knap-005-ux441ux443ux43fux435ux440ux432ux43eux437ux440ux430ux441ux442ux430ux44eux449ux438ux439-ux43bux451ux433ux43aux438ux439-ux440ux430ux43dux435ux446}

\hypertarget{swr-knap-005-супервозрастающий-лёгкий-ранец}{}

\hypertarget{swr-knap-005}{}

Для супервозрастающего варианта лёгкий ранец обязан строго возрастать: каждый элемент больше суммы всех предыдущих. Это свойство используется и жадной расшифровкой, и проверкой ответа. Последовательность, переданная или построенная с нарушением свойства, не образует допустимую криптосистему данного варианта.

\textbf{Требование.} Для супервозрастающего варианта лёгкий ранец должен быть приведён к строго сверхвозрастающей последовательности заданной размерности.

Приведение выполняется при генерации; размерность соответствует параметрам задачи. Проверка ответа повторно устанавливает свойство для переданного ранца --- см. {\hyperlink{swr-knap-009}{SWR-KNAP-009}}.

Источники: {\href{../04\_Системные\_требования/DOC-05\_Системные\_требования.pdf\#sr-knap-001}{SR-KNAP-001}}, {\href{../04\_Системные\_требования/DOC-05\_Системные\_требования.pdf\#sr-knap-004}{SR-KNAP-004}}.

\subsection{SWR-KNAP-006. Лёгкий ранец на степенях}\label{swr-knap-006-ux43bux451ux433ux43aux438ux439-ux440ux430ux43dux435ux446-ux43dux430-ux441ux442ux435ux43fux435ux43dux44fux445}

\hypertarget{swr-knap-006-лёгкий-ранец-на-степенях}{}

\hypertarget{swr-knap-006}{}

Вариант на степенях строит лёгкий ранец из последовательных степеней выбранного основания. Такая последовательность автоматически сверхвозрастает, что сохраняет применимость общей схемы проверки. Основание выбирается из фиксированного набора простых чисел.

\textbf{Требование.} Для варианта на степенях элементы лёгкого ранца должны строиться как степени выбранного основания.

Выбранное основание сохраняется в данных генерации. После разбора корректных десятичных строк проверяющий сопоставляет целые значения элементов со степенями этого основания. Системное основание --- {\href{../04\_Системные\_требования/DOC-05\_Системные\_требования.pdf\#sr-knap-006-вариант-на-степенях}{вариант на степенях}}.

Источники: {\href{../04\_Системные\_требования/DOC-05\_Системные\_требования.pdf\#sr-knap-001}{SR-KNAP-001}}, {\href{../04\_Системные\_требования/DOC-05\_Системные\_требования.pdf\#sr-knap-006}{SR-KNAP-006}}.

\subsection{SWR-KNAP-007. Сохранение сгенерированной задачи}\label{swr-knap-007-ux441ux43eux445ux440ux430ux43dux435ux43dux438ux435-ux441ux433ux435ux43dux435ux440ux438ux440ux43eux432ux430ux43dux43dux43eux439-ux437ux430ux434ux430ux447ux438}

\hypertarget{swr-knap-007-сохранение-сгенерированной-задачи}{}

\hypertarget{swr-knap-007}{}

Построенное условие должно стать данными создаваемой сессии, а не промежуточным значением в памяти генератора. Вместе с условием сохраняются вариант и принадлежность задачи. Отказ генерации или сохранения обрабатывается общим правилом полного набора: частичная сессия не остаётся активной.

\textbf{Требование.} Построенная задача должна сохраняться вместе с вариантом и принадлежностью, а её условие --- передаваться в состав создаваемой сессии до подтверждения операции создания.

Сохранённое условие затем используется при проверке без повторной генерации. Транзакционные границы операции создания определены {\href{DOC-06.2\_Сессии\_задач\_и\_планировщик.pdf\#swr-ses-005-сохранение-созданного-набора}{правилом сохранения набора}} домена сессий; ранцевой обработчик участвует в ней как источник условия.

Источники: {\href{../04\_Системные\_требования/DOC-05\_Системные\_требования.pdf\#sr-knap-001}{SR-KNAP-001}}, {\href{../04\_Системные\_требования/DOC-05\_Системные\_требования.pdf\#sr-knap-002}{SR-KNAP-002}}; {\href{DOC-06\_Программные\_требования.pdf\#swr-sys-003}{SWR-SYS-003}}, {\href{DOC-06.2\_Сессии\_задач\_и\_планировщик.pdf\#swr-ses-004}{SWR-SES-004}}.

\section{Проверка решения}\label{4-ux43fux440ux43eux432ux435ux440ux43aux430-ux440ux435ux448ux435ux43dux438ux44f}

\hypertarget{4-проверка-решения}{}

Решение передаётся в полях, перечисленных в {\href{DOC-06.2\_Сессии\_задач\_и\_планировщик.pdf\#ранцевая-задача}{контракте ранцевой задачи}}. Сначала компонент проверяет форму, например что список весов передан массивом десятичных строк. Затем он проверяет значения по {\hyperlink{условие-и-его-состав}{выданному условию}}. Вердикт означает результат математической проверки. Ошибка формы оставляет предметную попытку доступной до общего срока.

\subsection{SWR-KNAP-008. Предварительная сверка ответа}\label{swr-knap-008-ux43fux440ux435ux434ux432ux430ux440ux438ux442ux435ux43bux44cux43dux430ux44f-ux441ux432ux435ux440ux43aux430-ux43eux442ux432ux435ux442ux430}

\hypertarget{swr-knap-008-предварительная-сверка-ответа}{}

\hypertarget{swr-knap-008}{}

Для обоих вариантов используется одна {\href{DOC-06.2\_Сессии\_задач\_и\_планировщик.pdf\#ранцевая-задача}{схема решения}}. Она задаёт десять обязательных полей. Размерность передаётся положительным целым JSON-числом, вариант --- строкой из выданного условия, исходное и расшифрованное сообщения --- массивами логических значений. Для элементов ранцев и остальных целых величин установлена запись десятичными строками. Ни одно обязательное поле не принимает \texttt{null}.

\textbf{Требование.} Перед математической проверкой ПО должно сверить состав и типы полей со схемой решения и сравнить строку варианта с вариантом сохранённого условия. При несовпадении варианта или нарушении схемы клиент получает \texttt{400 INVALID\_\allowbreak{}REQUEST} без предметного вердикта.

Переданный вариант сравнивается с вариантом сохранённого условия: совпадение обязательно, подмена не допускается. Ошибка формы сохраняет попытку; далее допустимый ответ передаётся математической проверке. Разделение стадий следует {\href{DOC-06.2\_Сессии\_задач\_и\_планировщик.pdf\#swr-ses-009-проверка-формата}{правилу проверки формата}} домена сессий.

Источники: {\href{../04\_Системные\_требования/DOC-05\_Системные\_требования.pdf\#sr-knap-003}{SR-KNAP-003}}; {\href{../04\_Системные\_требования/DOC-05\_Системные\_требования.pdf\#sr-ses-009}{SR-SES-009}}, {\href{DOC-06.2\_Сессии\_задач\_и\_планировщик.pdf\#swr-ses-009}{SWR-SES-009}}, {\href{DOC-06.2\_Сессии\_задач\_и\_планировщик.pdf\#swr-ses-022}{SWR-SES-022}}.

\subsection{SWR-KNAP-009. Пошаговая проверка величин}\label{swr-knap-009-ux43fux43eux448ux430ux433ux43eux432ux430ux44f-ux43fux440ux43eux432ux435ux440ux43aux430-ux432ux435ux43bux438ux447ux438ux43d}

\hypertarget{swr-knap-009-пошаговая-проверка-величин}{}

\hypertarget{swr-knap-009}{}

Предметная проверка перебирает величины криптосистемы по установленному порядку: длину и элементы сообщений и ранцев, сверхвозрастаемость, соотношение модуля с суммой лёгкого ранца, взаимную простоту множителя, тяжёлый ранец, обратный множитель, шифртекст и расшифрованное сообщение. Расшифровка выполняется жадным алгоритмом по супервозрастающей последовательности; дополнительно устанавливается положительность элементов ранцев.

\textbf{Требование.} ПО должно проверять величины, перечисленные в {\href{../04\_Системные\_требования/DOC-05\_Системные\_требования.pdf\#sr-knap-004}{SR-KNAP-004}}, поэтапно, включая расшифровку жадным алгоритмом по супервозрастающей последовательности и положительность элементов ранцев.

Для обоих вариантов проверка охватывает расшифрованное сообщение. Если остальные величины верны, но расшифрованное сообщение отличается от ожидаемого, результат проверки отрицательный. Вычисление не запускает собственных фоновых процессов. Его длительность входит в {\href{DOC-06.2\_Сессии\_задач\_и\_планировщик.pdf\#swr-ses-020-предел-обработки}{предел обработки}} принятого ответа. При истечении предела операция завершается техническим отказом без предметного вердикта.

Источник: {\href{../04\_Системные\_требования/DOC-05\_Системные\_требования.pdf\#sr-knap-004}{SR-KNAP-004}}; {\href{../04\_Системные\_требования/DOC-05\_Системные\_требования.pdf\#sr-ses-012}{SR-SES-012}}.

\subsection{SWR-KNAP-010. Диагностика неверной величины}\label{swr-knap-010-ux434ux438ux430ux433ux43dux43eux441ux442ux438ux43aux430-ux43dux435ux432ux435ux440ux43dux43eux439-ux432ux435ux43bux438ux447ux438ux43dux44b}

\hypertarget{swr-knap-010-диагностика-неверной-величины}{}

\hypertarget{swr-knap-010}{}

Первая несовпавшая величина определяет текст диагностики. Пояснение относится к завершённой неверной попытке и позволяет студенту найти ошибку в собственной цепочке вычислений. Дальнейшие величины после первого несовпадения не сравниваются: вердикт уже определён.

\textbf{Требование.} При первом несовпадении ПО должно возвращать вердикт «неверно» с указанием величины, не прошедшей проверку.

После фиксации вердикта интерфейс возвращает \texttt{200\ OK}. Верный ответ даёт \texttt{isOk\ =\ true} и \texttt{errorMessage\ =\ null}. При неверном ответе \texttt{isOk\ =\ false}, а строка \texttt{errorMessage} указывает первую величину, для которой обнаружена ошибка. Оба вердикта завершают попытку. Категории отказов определены в {\href{DOC-06.2\_Сессии\_задач\_и\_планировщик.pdf\#коды-и-представление-ошибок}{контракте ошибок}}.

Источник: {\href{../04\_Системные\_требования/DOC-05\_Системные\_требования.pdf\#sr-knap-005}{SR-KNAP-005}}.

\subsection{SWR-KNAP-011. Защита итогового состояния}\label{swr-knap-011-ux437ux430ux449ux438ux442ux430-ux438ux442ux43eux433ux43eux432ux43eux433ux43e-ux441ux43eux441ux442ux43eux44fux43dux438ux44f}

\hypertarget{swr-knap-011-защита-итогового-состояния}{}

\hypertarget{swr-knap-011}{}

Повторное обращение к завершённой задаче не должно запускать математическую проверку. Решение принимается по сохранённым данным задачи, поэтому состояние клиента и повторные соединения не влияют на запрет второй попытки. Восстановление отображения выполняется чтением сессии. Если повторное прохождение контрольной разрешено, оно создаёт отдельную сессию с новыми задачами. Запрет второй попытки на прежней задаче продолжает действовать.

\textbf{Требование.} Если задача уже завершена, управление сессиями должно отклонить новый ответ до вызова ранцевого проверяющего. Сохранённый результат остаётся доступным для чтения.

Для владельца итоговой задачи код отказа зависит от времени получения полного запроса. До общего срока сессии действует \texttt{409 TASK\_\allowbreak{}FINALIZED}, начиная с общего срока --- \texttt{410 SESSION\_\allowbreak{}EXPIRED}. Запрос к чужой либо отсутствующей задаче получает \texttt{404 TASK\_\allowbreak{}NOT\_\allowbreak{}FOUND}. Порядок установлен {\href{DOC-06.2\_Сессии\_задач\_и\_планировщик.pdf\#swr-ses-008-проверка-допуска}{правилами допуска}} и {\href{DOC-06.2\_Сессии\_задач\_и\_планировщик.pdf\#swr-ses-013-сохранение-итогового-состояния}{сохранения итогового состояния}}.

Источники: {\href{../04\_Системные\_требования/DOC-05\_Системные\_требования.pdf\#sr-knap-007}{SR-KNAP-007}}; {\href{../04\_Системные\_требования/DOC-05\_Системные\_требования.pdf\#sr-ses-008}{SR-SES-008}}, {\href{DOC-06.2\_Сессии\_задач\_и\_планировщик.pdf\#swr-ses-013}{SWR-SES-013}}.

\section{Редакционные генерации}\label{5-ux440ux435ux434ux430ux43aux446ux438ux43eux43dux43dux44bux435-ux433ux435ux43dux435ux440ux430ux446ux438ux438}

\hypertarget{5-редакционные-генерации}{}

Редакционный пример --- отдельное учебное задание с полным эталонным решением для самоподготовки. Можно запросить пример по названию варианта или задать часть математических параметров. Выбор программы, решающей задачу указанного варианта, ниже назван диспетчеризацией.

\subsection{SWR-KNAP-012. Генерация по варианту}\label{swr-knap-012-ux433ux435ux43dux435ux440ux430ux446ux438ux44f-ux43fux43e-ux432ux430ux440ux438ux430ux43dux442ux443}

\hypertarget{swr-knap-012-генерация-по-варианту}{}

\hypertarget{swr-knap-012}{}

Редакционный запрос без дополнительных параметров выдаёт эталон выбранного варианта. Недостающие данные выбираются системой: основание --- из фиксированного набора простых чисел, остальные параметры --- по правилам варианта. Результат содержит полный набор величин и служит материалом для разбора решения.

\textbf{Требование.} Запрос выбранного варианта без дополнительных параметров должен строить эталон с полным решением; параметры, не указанные в запросе, выбираются ПО, в том числе основание --- из фиксированного набора простых чисел.

Редакционная генерация не связана с сессиями: она не проверяет допуски и не расходует попытки. Состав выдачи соответствует {\href{../04\_Системные\_требования/DOC-05\_Системные\_требования.pdf\#sr-knap-008-редакционная-задача-по-варианту}{эталонной задаче}}.

Источник: {\href{../04\_Системные\_требования/DOC-05\_Системные\_требования.pdf\#sr-knap-008}{SR-KNAP-008}}.

\subsection{SWR-KNAP-013. Генерация с параметрами}\label{swr-knap-013-ux433ux435ux43dux435ux440ux430ux446ux438ux44f-ux441-ux43fux430ux440ux430ux43cux435ux442ux440ux430ux43cux438}

\hypertarget{swr-knap-013-генерация-с-параметрами}{}

\hypertarget{swr-knap-013}{}

Параметры преподавателя делают эталон воспроизводимым в методических целях: размерность, вариант, сообщение, лёгкий ранец и множитель задаются запросом. Система проверяет математическую допустимость до построения; нарушение взаимной простоты множителя с модулем делает обратное преобразование невозможным и отклоняется.

\textbf{Требование.} Запрос с параметрами должен использовать переданные размерность, вариант, сообщение, лёгкий ранец и множитель; взаимная простота множителя с модулем должна проверяться до построения эталона.

Отклонение сопровождается диагностикой причины; частично построенный эталон не выдаётся. Отказ редакционной генерации не влияет на сессии и попытки.

Источник: {\href{../04\_Системные\_требования/DOC-05\_Системные\_требования.pdf\#sr-knap-009}{SR-KNAP-009}}.

\subsection{SWR-KNAP-014. Диспетчеризация по варианту}\label{swr-knap-014-ux434ux438ux441ux43fux435ux442ux447ux435ux440ux438ux437ux430ux446ux438ux44f-ux43fux43e-ux432ux430ux440ux438ux430ux43dux442ux443}

\hypertarget{swr-knap-014-диспетчеризация-по-варианту}{}

\hypertarget{swr-knap-014}{}

В редакционной генерации вариант указывает преподаватель. Один и тот же выбор действует при запросе эталона без параметров и при задании собственных параметров. Для запроса на степенях лёгкий ранец должен состоять из степеней выбранного основания. Одной сверхвозрастаемости элементов недостаточно, чтобы признать такой эталон верным.

\textbf{Требование.} Редакционный запрос должен направляться решателю соответствующего варианта.

Выбор выполняется по варианту запроса; результат диспетчеризации --- эталон, соответствующий указанному варианту по составу величин.

Источники: {\href{../04\_Системные\_требования/DOC-05\_Системные\_требования.pdf\#sr-knap-008}{SR-KNAP-008}}, {\href{../04\_Системные\_требования/DOC-05\_Системные\_требования.pdf\#sr-knap-009}{SR-KNAP-009}}.

\subsection{SWR-KNAP-015. Генерация без сохранения}\label{swr-knap-015-ux433ux435ux43dux435ux440ux430ux446ux438ux44f-ux431ux435ux437-ux441ux43eux445ux440ux430ux43dux435ux43dux438ux44f}

\hypertarget{swr-knap-015-генерация-без-сохранения}{}

\hypertarget{swr-knap-015}{}

Редакционные примеры не являются попытками и не должны попадать в хранилище сессий. Результат строится на лету и возвращается напрямую; повторный запрос выполняет генерацию заново.

\textbf{Требование.} Результаты редакционных генераций не должны записываться в хранилище.

Правило не затрагивает сохранение задач, создаваемых в сессиях по {\hyperlink{swr-knap-007}{SWR-KNAP-007}}.

Источник: {\href{../04\_Системные\_требования/DOC-05\_Системные\_требования.pdf\#sr-knap-010}{SR-KNAP-010}}.

\section{Хранение}\label{6-ux445ux440ux430ux43dux435ux43dux438ux435}

\hypertarget{6-хранение}{}

\subsection{SWR-KNAP-016. Сохранность условий}\label{swr-knap-016-ux441ux43eux445ux440ux430ux43dux43dux43eux441ux442ux44c-ux443ux441ux43bux43eux432ux438ux439}

\hypertarget{swr-knap-016-сохранность-условий}{}

\hypertarget{swr-knap-016}{}

Условия сгенерированных ранцевых задач должны переживать перезапуск приложения: вместе с вариантом и принадлежностью они образуют данные, по которым выполняется проверка и чтение результата. Списковые величины хранятся в текстовом представлении колонок и восстанавливаются без изменения значений.

\textbf{Требование.} Условия сгенерированных ранцевых задач должны сохраняться между запусками приложения вместе с вариантом и принадлежностью; списковые величины --- в текстовом представлении; сохранённые условия доступны для просмотра истории.

Сохранённые условия читаются историей задач без повторной генерации. Общие границы сохранения и восстановления определены {\href{DOC-06\_Программные\_требования.pdf\#swr-sys-008-сохранность-завершённых-данных}{правилом сохранности}} DOC-06.

Источники: {\href{../04\_Системные\_требования/DOC-05\_Системные\_требования.pdf\#sr-sys-007}{SR-SYS-007}}; {\href{DOC-06\_Программные\_требования.pdf\#swr-sys-008}{SWR-SYS-008}}.

\section{Трассировка и описательные проверки}\label{7-ux442ux440ux430ux441ux441ux438ux440ux43eux432ux43aux430-ux438-ux43eux43fux438ux441ux430ux442ux435ux43bux44cux43dux44bux435-ux43fux440ux43eux432ux435ux440ux43aux438}

\hypertarget{7-трассировка-и-описательные-проверки}{}

Таблица связывает каждое требование с системным основанием и наблюдаемым признаком его выполнения. В последней графе указаны ожидаемые результаты для разбора требований.

{\def\LTcaptype{none} % do not increment counter
\begin{longtable}[]{@{}
  >{\raggedright\arraybackslash}p{(\linewidth - 4\tabcolsep) * \real{0.2323}}
  >{\raggedright\arraybackslash}p{(\linewidth - 4\tabcolsep) * \real{0.2619}}
  >{\raggedright\arraybackslash}p{(\linewidth - 4\tabcolsep) * \real{0.5058}}@{}}
\toprule\noalign{}
\begin{minipage}[b]{\linewidth}\raggedright
SWR
\end{minipage} & \begin{minipage}[b]{\linewidth}\raggedright
Системные источники
\end{minipage} & \begin{minipage}[b]{\linewidth}\raggedright
Условия и ожидаемый результат проверки
\end{minipage} \\
\midrule\noalign{}
\endhead
\bottomrule\noalign{}
\endlastfoot
{\hyperlink{swr-knap-001}{SWR-KNAP-001}} & {\href{../04\_Системные\_требования/DOC-05\_Системные\_требования.pdf\#sr-knap-001}{SR-KNAP-001}} & Задачи, для которых шаблон указывает ранцевый квалификатор, создаются этим обработчиком; другие типы ему не передаются \\
{\hyperlink{swr-knap-002}{SWR-KNAP-002}} & {\href{../04\_Системные\_требования/DOC-05\_Системные\_требования.pdf\#sr-knap-001}{SR-KNAP-001}} & Разбор обоих исходов случайного выбора показывает соответствие сохранённого варианта выданному условию. Запрос студента не задаёт вариант \\
{\hyperlink{swr-knap-003}{SWR-KNAP-003}} & {\href{../04\_Системные\_требования/DOC-05\_Системные\_требования.pdf\#sr-knap-001}{SR-KNAP-001}}, {\href{../04\_Системные\_требования/DOC-05\_Системные\_требования.pdf\#sr-knap-002}{SR-KNAP-002}} & Длина сообщения от 3 до 5. При длине 3 хотя бы один элемент равен \texttt{false}, при длине 4 или 5 таких элементов не менее двух \\
{\hyperlink{swr-knap-004}{SWR-KNAP-004}} & {\href{../04\_Системные\_требования/DOC-05\_Системные\_требования.pdf\#sr-knap-002}{SR-KNAP-002}} & В условии скрыты все элементы ранца, кроме не более двух; множитель может отсутствовать \\
{\hyperlink{swr-knap-005}{SWR-KNAP-005}} & {\href{../04\_Системные\_требования/DOC-05\_Системные\_требования.pdf\#sr-knap-001}{SR-KNAP-001}}, {\href{../04\_Системные\_требования/DOC-05\_Системные\_требования.pdf\#sr-knap-004}{SR-KNAP-004}} & Лёгкий ранец супервозрастающего варианта выдерживает проверку сверхвозрастаемости \\
{\hyperlink{swr-knap-006}{SWR-KNAP-006}} & {\href{../04\_Системные\_требования/DOC-05\_Системные\_требования.pdf\#sr-knap-001}{SR-KNAP-001}}, {\href{../04\_Системные\_требования/DOC-05\_Системные\_требования.pdf\#sr-knap-006}{SR-KNAP-006}} & Элементы ранца варианта на степенях соответствуют степеням основания \\
{\hyperlink{swr-knap-007}{SWR-KNAP-007}} & {\href{../04\_Системные\_требования/DOC-05\_Системные\_требования.pdf\#sr-knap-001}{SR-KNAP-001}}, {\href{../04\_Системные\_требования/DOC-05\_Системные\_требования.pdf\#sr-knap-002}{SR-KNAP-002}} & Созданная сессия содержит сохранённые условия с вариантами и принадлежностью \\
{\hyperlink{swr-knap-008}{SWR-KNAP-008}} & {\href{../04\_Системные\_требования/DOC-05\_Системные\_требования.pdf\#sr-knap-003}{SR-KNAP-003}}, {\href{../04\_Системные\_требования/DOC-05\_Системные\_требования.pdf\#sr-ses-009}{SR-SES-009}} & Ответ с чужим вариантом или без обязательных полей отклоняется до вердикта \\
{\hyperlink{swr-knap-009}{SWR-KNAP-009}} & {\href{../04\_Системные\_требования/DOC-05\_Системные\_требования.pdf\#sr-knap-004}{SR-KNAP-004}}, {\href{../04\_Системные\_требования/DOC-05\_Системные\_требования.pdf\#sr-ses-012}{SR-SES-012}} & При неизменной допустимой форме искажённая величина даёт отрицательный вердикт. Если срок обработки истёк до фиксации, возвращается технический отказ \\
{\hyperlink{swr-knap-010}{SWR-KNAP-010}} & {\href{../04\_Системные\_требования/DOC-05\_Системные\_требования.pdf\#sr-knap-005}{SR-KNAP-005}} & Диагностика указывает первую несовпавшую величину \\
{\hyperlink{swr-knap-011}{SWR-KNAP-011}} & {\href{../04\_Системные\_требования/DOC-05\_Системные\_требования.pdf\#sr-knap-007}{SR-KNAP-007}}, {\href{../04\_Системные\_требования/DOC-05\_Системные\_требования.pdf\#sr-ses-008}{SR-SES-008}} & Повторный ответ после итогового состояния не создаёт второго вердикта \\
{\hyperlink{swr-knap-012}{SWR-KNAP-012}} & {\href{../04\_Системные\_требования/DOC-05\_Системные\_требования.pdf\#sr-knap-008}{SR-KNAP-008}} & Эталон содержит полный набор величин выбранного варианта \\
{\hyperlink{swr-knap-013}{SWR-KNAP-013}} & {\href{../04\_Системные\_требования/DOC-05\_Системные\_требования.pdf\#sr-knap-009}{SR-KNAP-009}} & Переданные параметры используются; невзаимно простой множитель --- отказ с причиной \\
{\hyperlink{swr-knap-014}{SWR-KNAP-014}} & {\href{../04\_Системные\_требования/DOC-05\_Системные\_требования.pdf\#sr-knap-008}{SR-KNAP-008}}, {\href{../04\_Системные\_требования/DOC-05\_Системные\_требования.pdf\#sr-knap-009}{SR-KNAP-009}} & По каждому из двух значений варианта возвращается эталон, удовлетворяющий математическому условию именно запрошенного варианта \\
{\hyperlink{swr-knap-015}{SWR-KNAP-015}} & {\href{../04\_Системные\_требования/DOC-05\_Системные\_требования.pdf\#sr-knap-010}{SR-KNAP-010}} & После редакционной генерации хранилище не пополняется \\
{\hyperlink{swr-knap-016}{SWR-KNAP-016}} & {\href{../04\_Системные\_требования/DOC-05\_Системные\_требования.pdf\#sr-sys-007}{SR-SYS-007}} & После перезапуска условия, варианты и принадлежность доступны \\
\end{longtable}
}

\section{Разбор граничных случаев}\label{8-ux440ux430ux437ux431ux43eux440-ux433ux440ux430ux43dux438ux447ux43dux44bux445-ux441ux43bux443ux447ux430ux435ux432}

\hypertarget{8-разбор-граничных-случаев}{}

Следующие случаи используют один исходный набор с каждой задачей в состоянии \texttt{NEW} (решение ещё не зафиксировано). Результаты \texttt{SOLVED} и \texttt{FAILED} и их влияние на набор определены в {\href{DOC-06.2\_Сессии\_задач\_и\_планировщик.pdf\#swr-ses-018-общая-функция-итога-сессии}{таблице состояний сессии}}. Отправки в таблице рассматриваются независимо. Кроме случая с истечением срока, запрос владельца принят своевременно и обработка успевает закончиться в пределах 60 секунд от получения запроса.

{\def\LTcaptype{none} % do not increment counter
\begin{longtable}[]{@{}
  >{\raggedright\arraybackslash}p{(\linewidth - 2\tabcolsep) * \real{0.3200}}
  >{\raggedright\arraybackslash}p{(\linewidth - 2\tabcolsep) * \real{0.6800}}@{}}
\toprule\noalign{}
\begin{minipage}[b]{\linewidth}\raggedright
Входные данные
\end{minipage} & \begin{minipage}[b]{\linewidth}\raggedright
Ожидаемый результат
\end{minipage} \\
\midrule\noalign{}
\endhead
\bottomrule\noalign{}
\endlastfoot
Клиент прислал полное решение, но указал в \texttt{type} другой вариант ранца & \texttt{400 INVALID\_\allowbreak{}REQUEST}. Обработчик не подставляет другой вариант, предметный результат отсутствует, попытка сохраняется \\
Обязательное поле \texttt{module} передано JSON-числом & \texttt{400 INVALID\_\allowbreak{}REQUEST}. Клиент может исправить тип и отправить новое решение до общего срока сессии \\
Вариант совпал, поля имеют заданные типы, но значение модуля математически неверно & При фиксации --- \texttt{200\ OK} с \texttt{isOk\ =\ false}; задача получает \texttt{FAILED}, а пояснение называет первую неверную величину \\
У задачи на степенях все величины верны, кроме одного бита расшифрованного сообщения & Отрицательный предметный результат указывает на расшифрованное сообщение. При чтении возвращается тот же вердикт \\
Каждая переданная величина верна & После фиксации --- \texttt{200\ OK}, задача в состоянии \texttt{SOLVED}, поле \texttt{errorMessage} присутствует со значением \texttt{null} \\
Проверяющий получил верный ответ, но вердикт не зафиксирован к моменту «получение полного запроса + 60 секунд» & \texttt{503 CHECK\_\allowbreak{}TIMEOUT}, предметный результат отсутствует. Повтор возможен только до общего срока открытой задачи \\
Владелец направил решение контрольной ранцевой задачи ровно в момент закрытия мероприятия, хотя с создания сессии ещё не прошло 20 минут & Подсистема сессий возвращает \texttt{410 SESSION\_\allowbreak{}EXPIRED} до вызова ранцевого обработчика. Своевременно полученный ответ может завершиться после закрытия в пределах исходных 60 секунд. \\
\end{longtable}
}

\section{История редакций}\label{9-ux438ux441ux442ux43eux440ux438ux44f-ux440ux435ux434ux430ux43aux446ux438ux439}

\hypertarget{9-история-редакций}{}

{\def\LTcaptype{none} % do not increment counter
\begin{longtable}[]{@{}
  >{\raggedright\arraybackslash}p{(\linewidth - 4\tabcolsep) * \real{0.2357}}
  >{\raggedright\arraybackslash}p{(\linewidth - 4\tabcolsep) * \real{0.2357}}
  >{\raggedright\arraybackslash}p{(\linewidth - 4\tabcolsep) * \real{0.5287}}@{}}
\toprule\noalign{}
\begin{minipage}[b]{\linewidth}\raggedright
Дата
\end{minipage} & \begin{minipage}[b]{\linewidth}\raggedright
Версия
\end{minipage} & \begin{minipage}[b]{\linewidth}\raggedright
Изменение
\end{minipage} \\
\midrule\noalign{}
\endhead
\bottomrule\noalign{}
\endlastfoot
07.11.2026 & 0.8 & Подготовлены программные требования к ранцевым задачам, связи с DOC-05 и описания проверок \\
20.11.2026 & 0.9 & Уточнены выбор варианта в сессии и в редакционной генерации, различие ошибок формы и вычисления, состав вердикта и допуск повторного обращения. Добавлены граничные случаи \\
\end{longtable}
}


\end{document}
